% ─────────────────────────────────────────────────────────────────
\section{Introduction}
\label{sec:intro}
% ─────────────────────────────────────────────────────────────────

Protein fold classification is a central problem in structural bioinformatics.
The three-dimensional surface shape of a protein, together with its
electrostatic potential field, encodes information about its biochemical
function and evolutionary family.
The SHREC\,25 Shape Retrieval Contest~\cite{shrec25} provides a benchmark that
casts fold recognition as a closed-world, 99-class classification task on
protein surface meshes, with a strongly class-imbalanced training set of
9\,244 proteins.

\paragraph{Why this problem is hard.}
Three structural challenges make SHREC\,25 particularly demanding:
\begin{enumerate}
  \item \textbf{Class imbalance.}
        The most frequent class (class\,8) has 2\,054 training examples while
        the rarest class has only 2 --- a 1\,027$\times$ ratio.
        Nine classes have zero test proteins; evaluating macro-averaged F1 over
        88 testable classes directly penalises poor recall on rare classes.

  \item \textbf{Near-inseparable class pairs.}
        Several protein families share highly similar surface geometry and
        electrostatics.
        The pair (class\,61, class\,90), for instance, cannot be reliably
        separated even with full-resolution VTK meshes and protein-level
        support-vector-machine classifiers (balanced accuracy $\approx 0.56$,
        barely above chance).

  \item \textbf{Surface-only feature budget.}
        The dataset provides only vertex positions, surface normals, and two
        scalar fields --- electrostatic potential (\textit{Potential}) and
        normal component of potential (\textit{NormalPotential}).
        Sequence information, volumetric electron density, and other richer
        representations are absent.
\end{enumerate}

\paragraph{Contributions.}
This report documents a complete end-to-end study with the following
contributions:
\begin{itemize}
  \item A preprocessing pipeline that converts VTK surface meshes into
        graph-learning-ready SSL graphs via farthest-point subsampling and
        k-nearest-neighbour edge construction
        (Section~\ref{sec:pipeline}).

  \item A statistical separability analysis at vertex level (Kolmogorov--Smirnov,
        Sliced Wasserstein), protein level (SVM on global statistics), and
        spatial level (Moran's~I autocorrelation) to characterise which
        features actually discriminate between classes
        (Section~\ref{sec:analysis}).

  \item A hierarchical graph neural network (\textit{HierGIN}) that stacks a
        node-level GIN over the 1\,000-node surface graph, a geodesic Voronoi
        patch pooling layer, and a patch-level GIN before global readout and
        classification (Section~\ref{sec:arch}).

  \item Systematic ablation of eight node-backbone types, five feature
        configurations, patch counts, hidden dimensions, and training
        regularisation choices (Section~\ref{sec:ablation}).

  \item A comparison with external methods (PointMamba~\cite{pointmamba},
        Mamba3D~\cite{mamba3d}) submitted to the SHREC\,25 contest, and an
        analysis of the theoretical performance ceiling imposed by the dataset
        structure (Section~\ref{sec:results}).
\end{itemize}

Our best model, \textbf{HierGIN-8D}, achieves a macro-F1 of \textbf{0.7983}
on 88 evaluated test classes, surpassing flat-GNN baselines by $+$15.8\,pp,
while the fundamental confusion between class pairs (61,\,90) and (70,\,56)
persists across all models and feature configurations tested.
